\documentclass[11pt,a4paper,reqno]{amsart}
\usepackage[utf8]{inputenc}
\usepackage[T1]{fontenc}
\usepackage[ngerman]{babel}
\usepackage{amsmath,amssymb,amsthm,amsfonts,mathtools}
\usepackage{geometry}
\usepackage{listings}
\usepackage{xcolor}
\usepackage{hyperref}

\geometry{a4paper,margin=2.5cm,marginparwidth=1.5cm,marginparsep=0.3cm}

% Theorem-Umgebungen
\newtheorem{theorem}{Theorem}[section]
\newtheorem{lemma}[theorem]{Lemma}
\newtheorem{proposition}[theorem]{Proposition}
\newtheorem{corollary}[theorem]{Corollary}

\theoremstyle{definition}
\newtheorem{definition}[theorem]{Definition}
\newtheorem{axiom}[theorem]{Axiom}
\newtheorem{example}[theorem]{Example}

\theoremstyle{remark}
\newtheorem{remark}[theorem]{Remark}

% Code-Listing Einstellungen für SageMath / Python
\definecolor{codegreen}{rgb}{0,0.6,0}
\definecolor{codegray}{rgb}{0.5,0.5,0.5}
\definecolor{codepurple}{rgb}{0.58,0,0.82}
\definecolor{backcolour}{rgb}{0.97,0.97,0.97}

\lstdefinestyle{sagestyle}{
    backgroundcolor=\color{backcolour},   
    commentstyle=\color{codegreen},
    keywordstyle=\color{magenta},
    numberstyle=\tiny\color{codegray},
    stringstyle=\color{codepurple},
    basicstyle=\ttfamily\footnotesize,
    breakatwhitespace=false,         
    breaklines=true,                 
    captionpos=b,                    
    keepspaces=true,                 
    numbers=left,                    
    numbersep=5pt,                  
    showspaces=false,                
    showstringspaces=false,
    showtabs=false,
    tabsize=4,
    upquote=true,
    literate=
      {ä}{{\"a}}1 {ö}{{\"o}}1 {ü}{{\"u}}1
      {Ä}{{\"A}}1 {Ö}{{\"O}}1 {Ü}{{\"U}}1
      {ß}{{\ss}}1 {≡}{{$\equiv$}}1
}
\lstset{style=sagestyle}

\hypersetup{
    unicode=true,
    colorlinks=true,
    linkcolor=blue,
    citecolor=blue,
    urlcolor=blue,
    pdftitle={Typentrennung, Projektionsäquivalenz und Faserstrukturen},
    pdfauthor={Thomas Hoffbauer}
}

% Mathematische Abkürzungen
\newcommand{\M}{\mathcal{M}}
\newcommand{\N}{\mathbb{N}}
\newcommand{\Pset}{\mathbb{P}}
\newcommand{\Z}{\mathbb{Z}}
\newcommand{\Q}{\mathbb{Q}}
\newcommand{\R}{\mathbb{R}}
\newcommand{\Fact}{\operatorname{Fact}}
\newcommand{\status}{\operatorname{status}}

\title[Typentrennung, Faserstruktur und SageMath-Verifikation]{Typentrennung, Projektionsäquivalenz und Faserstrukturen in semiprimen Modellräumen mit numerischer SageMath-Verifikation}

\author{Thomas Hoffbauer}
\address{Computer Institut Bamberg, Deutschland}
\email{research@cib.de}

\date{5. August 2026}

\begin{document}
\raggedbottom
\emergencystretch=2em

\begin{abstract}
Aufbauend auf der algebraischen Absicherung von Modellräumen $\M$ gegenüber dem Hauptsatz der Arithmetik formulieren wir in dieser Arbeit die exakte Typentrennung für zusammengesetzte Modellrepräsentanten. Wir zeigen, dass ein Tupel $(a,b;c)_\M$ im Allgemeinen kein Primzahltripel in $\N^3$ bildet, sondern ein strukturiertes Objekt der Produktmenge $\Pset^2 \times \M$ darstellt. Durch Einführung der Projektionsäquivalenz $x \equiv_\pi y \iff \pi(x) = \pi(y)$ wird die Modellbeziehung $a \star b \star E_1 \equiv_\pi c \star E_2$ präzisiert, ohne Gleichheit oder Assoziiertheit der Zustände im Modellraum zu erzwingen. Wir definieren die Faserstruktur $M_n = \pi^{-1}(n) \subseteq \M$ und entwerfen eine dreistufige Datenstruktur (Arithmetische Ebene, Zustandsebene, Brückenebene) zur konsistenten Repräsentation semiprimer Invarianten. Ergänzend stellen wir ein vollständiges SageMath 10.5 Skript bereit, das die theoretischen Resultate am konkreten Modellraum $\Z[i]$ numerisch verifiziert.
\end{abstract}

\maketitle

\section{Einleitung}

In vorangegangenen Untersuchungen \cite{Hoffbauer2026} wurde gezeigt, dass die Koexistenz einer semiprimen Faktordarstellung $x = a \star b \star E_1$ und einer kompakten Zustandsdarstellung $y = c \star E_2$ im Modellraum $\M$ bei identischer Skalarprojektion $\pi(x) = \pi(y) = n \in \N$ keinen Widerspruch zum Hauptsatz der Arithmetik darstellt.

Die vorliegende Arbeit präzisiert die daraus folgenden formalen Konsequenzen. Insbesondere wird gezeigt, dass das Notationsobjekt $(a,b,c)$ einer strengen Typentrennung bedarf: Während $a$ und $b$ klassische Primfaktoren in $\N$ sind, ist $c$ ein Element des Modellraums $\M$. Wir analysieren die Faserstruktur $M_n = \pi^{-1}(n)$, stellen die Projektionsäquivalenz $\equiv_\pi$ auf ein rigides axiomatisches Fundament und implementieren eine SageMath-Referenzarchitektur zur auditierbaren Verifikation der Faser-Eigenschaften.

\section{Typentrennung und Kategoriale Einordnung}

\begin{definition}[Semiprimes Modell-Tupel]
Es sei $\Pset \subset \N$ die Menge der natürlichen Primzahlen und $\M$ ein Modellraum mit Produktoperation $\star: \M \times \M \to \M$ sowie Sektoroperatoren $E_1, E_2 \in \M$.
Ein \emph{semiprimes Modell-Tupel} ist ein Element
\begin{equation}
(a,b \,; c)_\M \;\in\; \Pset^2 \times \M.
\end{equation}
Das Semikolon trennt hierbei explizit die skalaren Arithmetik-Bausteine $(a,b) \in \Pset^2$ vom Modellzustand $c \in \M$.
\end{definition}

\begin{proposition}[Ausschlussregel für Standard-Primtriplets]
Es sei $\pi: \M \to \N$ eine Projektion mit $\pi(E_2) = 1$ und $\pi(c \star E_2) = \pi(c)$. Angenommen, $c$ wäre eine gewöhnliche Primzahl $c \in \Pset \subset \N$ mit kanonischer Einbettung $\pi(c) = c$.
Falls für echte Primzahlen $a,b \in \Pset$ mit $a,b > 1$ die Skalargleichheit $\pi(a \star b \star E_1) = \pi(c \star E_2)$ gilt und die Projektion auf der Faktorseite streng multiplikativ ist ($\pi(a \star b \star E_1) = a \cdot b$), so ist die Annahme $c \in \Pset$ falsch.
\end{proposition}

\begin{proof}
Aus den Prämissen folgt $a \cdot b = \pi(c \star E_2) = \pi(c) = c$. Da $a,b > 1$ echte Primzahlen sind, ist das Produkt $n = a \cdot b$ in $\N$ zusammengesetzt. Eine natürliche Zahl $c \in \N$ mit $c = a \cdot b$ besitzt mindestens zwei echte Teiler $a$ und $b$ und erfüllt $c \notin \Pset$.
\end{proof}

\begin{corollary}[Bedingungskatalog für kompakte Darstellungen]
\label{cor:ausschluss}
Für jedes Modell-Tupel $(a,b \,; c)_\M$ mit $\pi(a \star b \star E_1) = \pi(c \star E_2) = a \cdot b$ muss mindestens eine der folgenden Aussagen zutreffen:
\begin{enumerate}
    \item $c \notin \N$, sondern $c \in \M \setminus \N$ ist ein reiner Modellzustand.
    \item $c$ ist irreduzibel in $\M$, aber sein Skalarbild $\pi(c)$ ist zusammengesetzt in $\N$.
    \item Die Abbildung $\pi$ ist nicht-identisch auf $c$, d.\,h. $\pi(c) \neq c$.
    \item Der Sektoroperator $E_2$ ist nicht projektionsneutral ($\pi(E_2) \neq 1$).
    \item Die Projektionsabbildung $\pi$ ist nicht streng multiplikativ auf $a \star b \star E_1$.
\end{enumerate}
\end{corollary}

\section{Projektionsäquivalenz und Faserstrukturen}

\begin{definition}[Projektionsäquivalenz]
Auf dem Modellraum $\M$ definieren wir die zweistellige Relation der \emph{Projektionsäquivalenz} $\equiv_\pi \;\subseteq \M \times \M$ durch:
\begin{equation}
x \equiv_\pi y \quad :\Longleftrightarrow \quad \pi(x) = \pi(y).
\end{equation}
\end{definition}

\begin{lemma}[Eigenschaften der Projektionsäquivalenz]
Die Relation $\equiv_\pi$ ist eine Äquivalenzrelation auf $\M$. Die Äquivalenzklassen $[x]_\pi$ stehen in bijektiver Korrespondenz zu den Fasern $M_n$ der Projektion $\pi$:
\begin{equation}
M_n := \pi^{-1}(n) = \{ z \in \M \mid \pi(z) = n \}, \qquad n \in \operatorname{Im}(\pi) \subseteq \N.
\end{equation}
\end{lemma}

\begin{theorem}[Faser-Zerlegung und Invarianten-Erhaltung]
\label{thm:faser}
Der Modellraum $\M$ lässt sich als disjunkte Vereinigung seiner Skalarfasern darstellen:
\begin{equation}
\M = \bigsqcup_{n \in \operatorname{Im}(\pi)} M_n.
\end{equation}
Für zwei Modellzustände $x = a \star b \star E_1$ und $y = c \star E_2$ gilt:
\begin{equation}
x \equiv_\pi y \iff x,y \in M_n \quad \text{mit } n = a \cdot b.
\end{equation}
Aus $x \equiv_\pi y$ folgt weder $x = y$ noch $x \sim y$ (Assoziiertheit in $\M$), noch $\Fact(x) = \Fact(y)$.
\end{theorem}

\begin{proof}
Die Zerlegung folgt direkt aus den Eigenschaften von Urbildfasern einer mathematischen Abbildung. Dass $x \equiv_\pi y$ weder Gleichheit noch Assoziiertheit impliziert, wurde in \cite{Hoffbauer2026} durch das Gegenbeispiel in den Gaußschen Zahlen $\Z[i]$ mit $x=5$ und $y=3+4i$ bei Norm $\pi(z)=N(z)=25$ bewiesen.
\end{proof}

\section{Bedingungsanalysen für Multiplikativität}

In der Modellbildung muss präzise unterschieden werden, ob die Projektion $\pi$ global multiplikativ wirkt oder nur lokal Skalarwerte abgleicht.

\begin{proposition}[Multiplikativer Brückenschluss]
Es sei $\pi: \M \to \N$ eine streng multiplikative Abbildung mit $\pi(u \star v) = \pi(u) \cdot \pi(v)$ für alle $u,v \in \M$. Ferner gelte $\pi(a) = a$, $\pi(b) = b$ für $a,b \in \Pset$ sowie $\pi(E_1) = \pi(E_2) = 1$.
Dann gilt für $x = a \star b \star E_1$ und $y = c \star E_2$:
\begin{equation}
x \equiv_\pi y \implies \pi(c) = a \cdot b.
\end{equation}
\end{proposition}

\begin{proof}
Aufgrund der Multiplikativität gilt:
\begin{equation}
\pi(x) = \pi(a \star b \star E_1) = \pi(a) \cdot \pi(b) \cdot \pi(E_1) = a \cdot b \cdot 1 = a \cdot b.
\end{equation}
Ebenso gilt $\pi(y) = \pi(c \star E_2) = \pi(c) \cdot \pi(E_2) = \pi(c) \cdot 1 = \pi(c)$.
Aus $x \equiv_\pi y$ folgt $\pi(x) = \pi(y)$, mithin $\pi(c) = a \cdot b$.
\end{proof}

\begin{remark}[Status-Dichotomie für das kompakte Element $c$]
Unter den Voraussetzungen der obigen Proposition besitzt das Element $c \in \M$ eine zweifache Charakterisierung:
\begin{enumerate}
    \item Im Modellraum: $\status_\M(c) = \text{irreduzibel}$.
    \item Im Skalarraum: $\status_\N(\pi(c)) = \text{zusammengesetzt}$ (da $\pi(c) = a \cdot b$).
\end{enumerate}
Dies demonstriert, dass Begriffsbildungen wie ``irreduzibel'' strukturspezifisch sind und nicht ungeprüft über Projektionsgrenzen hinweg übertragen werden dürfen.
\end{remark}

\section{Dreistufige Formale Datenstruktur}

Zur konsistenten rechnergestützten und formalen Repräsentation semiprimer Invarianten im Modellraum $\M$ definieren wir ein dreistufiges Datenmodell.

\begin{definition}[Modell-Datensatz]
Ein \emph{semiprimer Modell-Datensatz} $\mathcal{R}_{a,b}$ ist ein 3-Tupel
\begin{equation}
\mathcal{R}_{a,b} := \left( \mathcal{A}(n), \, \Fact_\N(n), \, \mathcal{M}_n^{\operatorname{rep}} \right),
\end{equation}
wobei die Komponenten wie folgt definiert sind:
\begin{enumerate}
    \item \textbf{Arithmetische Ebene $\mathcal{A}(n)$:} Der skalare Hauptwert $n = a \cdot b \in \N$.
    \item \textbf{Primfaktorebene $\Fact_\N(n)$:} Die eindeutige arithmetische Bewertungsdarstellung:
    \begin{equation}
    \Fact_\N(n) = \{ (p, v_p(n)) \mid p \in \Pset, \, v_p(n) > 0 \} = \begin{cases} \{ (a,1), (b,1) \}, & a \neq b, \\ \{ (a,2) \}, & a = b. \end{cases}
    \end{equation}
    \item \textbf{Brücken- und Zustandsebene $\mathcal{M}_n^{\operatorname{rep}}$:} Eine gewählte Teilmenge repräsentativer Modellzustände aus der Faser $M_n$:
    \begin{equation}
    \mathcal{M}_n^{\operatorname{rep}} = \{ a \star b \star E_1, \; c \star E_2 \} \;\subseteq\; M_n.
    \end{equation}
\end{enumerate}
\end{definition}

Durch diese Architektur wird die Vermischung von skalaren Invarianten und geometrisch-algebraischen Modellrepräsentanten auf Datenebene unmöglich gemacht.

\section{Algorithmetische Referenzimplementierung in SageMath 10.5}

Zur auditierbaren Überprüfung der theoretischen Eigenschaften der Faserstrukturen $M_n = \pi^{-1}(n)$ und der Projektionsäquivalenz $\equiv_\pi$ dient der nachfolgende SageMath-Quellcode. Als konkreter Modellraum wird der Gitterring der Gaußschen Zahlen $\M = \Z[i]$ mit der algebraischen Norm $N(z) = z \bar{z}$ als Projektion $\pi$ genutzt.

\begin{lstlisting}[language=Python, caption={SageMath 10.5 Verifikationsscript für Fasern und Äquivalenzen}]
from sage.all import *

class SemiprimeModelSpace:
    def __init__(self, R):
        """
        Initialisiert den Modellraum M über einem Basisring R (z.B. Gauss'sche Zahlen ZZ[i]).
        """
        self.R = R
        self.P = Primes()

    def projection(self, element):
        """
        Projektionsabbildung pi: M -> N.
        Hier beispielhaft als algebraische Norm N(z) definiert.
        """
        return ZZ(element.norm())

    def compute_fiber(self, n, search_bound=20):
        """
        Berechnet die Faser M_n = pi^{-1}(n) innerhalb einer Suchgrenze.
        """
        fiber = []
        for x in range(-search_bound, search_bound + 1):
            for y in range(-search_bound, search_bound + 1):
                elem = self.R(x + y*I)
                if elem != 0 and self.projection(elem) == n:
                    fiber.append(elem)
        return fiber

    def is_irreducible_in_M(self, element):
        """
        Prüft, ob ein Element im Modellraum M irreduzibel ist.
        """
        return element.is_irreducible()

    def create_semiprime_dataset(self, a, b, E1=1, E2=1, c_candidate=None):
        """
        Erstellt das dreistufige Datenmodell R_{a,b} = (A(n), Fact_N(n), M_n^rep).
        """
        assert a in self.P and b in self.P, "a und b müssen arithmetische Primzahlen sein!"
        
        # 1. Arithmetische Ebene A(n)
        n = ZZ(a * b)
        
        # 2. Primfaktorebene Fact_N(n)
        fact_N = factor(n)
        
        # Modellzustand 1: x = a * b * E1 (Semiprim-Faktordarstellung)
        x = self.R(a * b * E1)
        
        # Modellzustand 2: y = c * E2 (Kompakter Zustand)
        fiber_n = self.compute_fiber(n)
        if c_candidate is None:
            irreducibles_in_fiber = [elem for elem in fiber_n if self.is_irreducible_in_M(elem)]
            c = irreducibles_in_fiber[0] if irreducibles_in_fiber else None
        else:
            c = self.R(c_candidate)
            
        y = c * self.R(E2) if c is not None else None
        
        # 3. Brücken- und Zustandsebene (Projektionsäquivalenzprüfung)
        pi_x = self.projection(x)
        pi_y = self.projection(y) if y is not None else None
        equiv_pi = (pi_x == pi_y) if y is not None else False

        dataset = {
            "Arithmetische_Ebene_n": n,
            "Primfaktorebene_Fact_N": fact_N,
            "Modellzustand_x": x,
            "Modellzustand_y": y,
            "c_is_irreducible_in_M": self.is_irreducible_in_M(c) if c else None,
            "pi(x)": pi_x,
            "pi(y)": pi_y,
            "Projektionsaequivalenz_x_equiv_pi_y": equiv_pi,
            "Faser_M_n_Groesse": len(fiber_n),
            "Faser_M_n_Elemente": fiber_n
        }
        return dataset

# ---------------------------------------------------------
# Test & Verifikation
# ---------------------------------------------------------
if __name__ == "__main__":
    M_Gauss = SemiprimeModelSpace(ZZ[I])
    a_prime = 5
    b_prime = 5

    ds = M_Gauss.create_semiprime_dataset(a=a_prime, b=b_prime, E1=1, E2=1, c_candidate=3+4*I)

    print("="*60)
    print(f"SEMIPRIMER MODELL-DATENSATZ FÜR n = {ds['Arithmetische_Ebene_n']}")
    print("="*60)
    print(f"1. Skalarwert n         : {ds['Arithmetische_Ebene_n']}")
    print(f"2. Arithmetische Fakten : {ds['Primfaktorebene_Fact_N']}")
    print(f"3. Modellzustand x      : {ds['Modellzustand_x']} (pi(x) = {ds['pi(x)']})")
    print(f"4. Modellzustand y      : {ds['Modellzustand_y']} (pi(y) = {ds['pi(y)']})")
    print(f"   -> c irreduzibel in M: {ds['c_is_irreducible_in_M']}")
    print(f"5. x ≡_pi y (Skalar)    : {ds['Projektionsaequivalenz_x_equiv_pi_y']}")
    print(f"6. Identisch (x == y)?  : {ds['Modellzustand_x'] == ds['Modellzustand_y']}")
    print(f"7. Faser M_{ds['Arithmetische_Ebene_n']} Größe : {ds['Faser_M_n_Groesse']} Elemente")
    print(f"   Faser M_{ds['Arithmetische_Ebene_n']} Auszug: {ds['Faser_M_n_Elemente'][:8]}...")
    print("="*60)
\end{lstlisting}

\section{Grenzen der Konsistenzarchitektur}

\begin{remark}[Nicht-bewiesene Aussagen / Offene Fragen]
Das vorliegende Framework stellt eine reine Konsistenz- und Schutzarchitektur bereit. Folgende Existenz- und Klassifikationsaussagen sind \textbf{nicht} Gegenstand dieser Theorie und müssen für konkrete Ausprägungen von $\M$ separat bewiesen werden:
\begin{enumerate}
    \item \textbf{Existenz:} $\forall a,b \in \Pset \; \exists c \in \M: c \text{ irreduzibel in } \M \text{ und } \pi(c) = a \cdot b$.
    \item \textbf{Eindeutigkeit:} Die Eindeutigkeit des kompakten Repräsentanten $c$ innerhalb einer Faser $M_n$ bis auf Einheiten.
    \item \textbf{Konstruktiver Algorithmus:} Die explizite Abbildung oder Vorschrift $(a,b) \mapsto c$.
\end{enumerate}
\end{remark}

\section{Fazit}

Die Differenzierung zwischen dem skalaren Invariantenwert $n \in \N$ und den Elementen der Modellfaser $M_n = \pi^{-1}(n)$ löst sämtliche Scheinwidersprüche bezüglich des Hauptsatzes der Arithmetik auf. Das semiprime Objekt $(a,b \,; c)_\M$ ist kein Primzahltripel, sondern beschreibt die Äquivalenz zweier unterschiedlicher Modellzustände $a \star b \star E_1 \equiv_\pi c \star E_2$ bezüglich ihrer Projektion. Die Gültigkeit der eindeutigen Primfaktorzerlegung in $\N$ bleibt von der geometrischen oder algebraischen Vielfalt innerhalb der Fasern $M_n$ vollständig unberührt und lässt sich über die vorgestellte SageMath-Architektur präzise verifizieren.

\begin{thebibliography}{99}
\bibitem{Hoffbauer2026} T. Hoffbauer, \emph{On the Compatibility of Model-Space Representations with the Fundamental Theorem of Arithmetic}, arXiv preprint, August 2026.
\bibitem{Niven} I. Niven, H. S. Zuckerman, H. L. Montgomery, \emph{An Introduction to the Theory of Numbers}, 5th ed., John Wiley \& Sons, 1991.
\bibitem{Conway} J. H. Conway, D. A. Smith, \emph{On Quaternions and Octonions: Their Geometry, Arithmetic, and Symmetry}, A K Peters/CRC Press, 2003.
\end{thebibliography}

\end{document}